\section{Karma and Reincarnation}

\begin{quotex}
When you speak of your memory of the future, you're probably thinking of transmigration or reincarnation. Maybe you think I can help you fulfill your desinty. But the idea of reincarnation isn't mentioned anywhere in the old texts. Rather it emerged from the detritus of the Flood and is linked to the primeval serpent and to the dark people of prehistoric times.
\flright{\textsc{Miguel Serrano}, \textit{EL/ELLA}}
\end{quotex}


In his \emph{Maschera e volto dello spiritualismo contemporaneo}, Julius Evola offers a critique of Theosophy focusing on the notions of Karma and Reincarnation. That text is now available in English in the library section of this web site\footnote{\url{https://gornahoor.net/library}}.

Within the confines of a single chapter, Evola recapitulates the arguments of Rene Guenon against Theosophy in general and reincarnation in particular, and adds some of his own twists. It is a model of a fair and incisive criticism. It is fair because it represents the Theosophical position accurately; a Theosophist will likely agree with it even while coming to an opposite valuation. It is incisive because Evola exposes the unstated root assumptions behind the notion and then draws out the logical consequences.

Evola begins by contrasting the traditional notion of karma with the modern version as represented in Theosophy. In the Traditional conception, karma is the law of cause and effect, though it extends well beyond merely physical phenomena. The objection to reincarnation is both subtle and profound. It cannot be emphasized sufficiently that this is not an optional point of view but is fundamental to Traditionalism, as all its parts are intimately related into a coherent whole. I have met \emph{soi-disant} ``Traditionalists'', even among those who deign to publish books and journals, who reject what Evola and Guenon say about reincarnation.

Besides the main thread of the argument, which must be followed closely, note also the little digressions that are also of interest. Pay attention to what he writes about ``Providence'', since it will come up again in the second part of The Individual and the Becoming of the World. He also relates the Traditional teachings to some Catholic dogmas, such as the Fall and purgatory. Note how he cuts through the dogma and moralism to get to the core of the teachings.


\flrightit{Posted on 2008-09-03 by Cologero}

\begin{center}* * *\end{center}

\begin{footnotesize}
\begin{sffamily}
\texttt{James O'Meara on 2010-10-24 at 12:54 said:}

Cologero,

I was looking for a cheap copy of Meditations on the Tarot since you seem to think highly of it; looking at it on Amazon, it seems Tonberg is pro-reincarnation, viewing it as something the Hermeticist ``just knows from experience'' and disparaging arguments for or against. Can you reconcile T. and E. on this?

\hfill

\texttt{Cologero on 2010-10-24 at 13:53 said:}

You can find a PDF copy of the book, as well as the Marseille deck, in the Gornahoor library. It is an unusually clear presentation of Hermetism, though avoid the cults that are growing up around Tomberg. He is in the tradition of Eliphas Levi. I am sure you are aware that Aleister Crowley regarded himself as a reincarnation of Levi. (Crowley's notes to Levi's Keys to the Mysteries are illuminating.)

I am working on a larger piece on this very topic, but I'll provide a preview here. First of all, Tomberg experiences reincarnation as a ``psychic fact''. Of course, you cannot argue against a fact, yet a soul (the psychic level) does not incarnate. This needs to be understood in the context of what Tomberg writes about mythological archetypes repeating in time and what Guenon writes about psychic residues.

\hfill

\texttt{James O'Meara on 2010-10-24 at 18:30 said:}

I hadn't noticed that, thanks for making it available. Although I'm an old fashioned sort and will likely still get a book version at some point. Looking forward to your ``larger piece''.

\hfill

\texttt{Boreas on 2011-12-15 at 06:57 said:}

The theosophical teachings about karma and reincarnation are certainly many times confusing and overtly moralising, together with the ``evolutionary'' and humanistic garb, and I'm more inclined towards the Evolian and Guénonian viewpoint myself also.

Yet, I do think that Evola mis-interpreted the teachings about `elder and lesser brothers'* and the spiritual concept of `brotherhood' also, which is (or should be) nothing but the practical application of buddhi or spiritual unity, together with basic empathy among men and other beings. It doesn't necessarily presuppose levelling egalitarianism etc. at all, although it does make power-hungry nihilism and oppression of other people on the basis of `might is right' impossible.

* the concept of the transmigration of the ``spiritual monad'' can be traced at least as far as Plato, and I think this is the concept behind the theosophical teaching about `elder' and `lesser' brothers.

It is probable that he did this because of theosophical misunderstandigns and confusions, but I think there would be a real need for a re-conciliation between theosophists and traditionalists in this issue; both Evola and Guénon did believe in some sort of reincarnation themselves also. Terminological confusions should not be or become a separating wall between men and women who strive for a common goal (= spiritual elevation and evolution).

\hfill

\texttt{Nick on 2022-02-02 at 19:43 said:}

This is an interesting topic. The two comments I have to make on my part are the following.

First, Guenon writes more about what reincarnation is not rather than what it is. Some of the formulations he makes about the subject are strange. It is not exactly clear that he completely denies the teaching of reincarnation, there are some enigmatic points in his work about the subject. Rather it shows that Guenon wanted one more to take a distance from the way theosophy and spiritualism perceived these concepts and the most dominant representations.

Secondly, reincarnation is a canonical traditional teaching (even in initiatic organizations that Guenon considered completely canonical and the most esoteric --- Guenon is wrong on this point when he says that these teachings are more symbolical or allegorical in these traditions) a teaching from ancient schools to the present day and is not simply ``a modern teaching that is an invention of Theosophy''.

\hfill

\texttt{Noct on 2023-09-02 at 15:04 said:}

Traditionally understood reincarnation refers to a supra-individual principle (the transcendental Self, the Nous, the ``image and likeness'' or Spirit), which could be formulated as a ``modality'' or ``idea'' of God, whose ``project'' of ``realization'' ``implies passing through successive incarnations. The error of modern reincarnation is to conceive the focus or subject of the process to the ``conditioned self'', the psyche or the soul, without recognizing that this is in fact the ``wrapping'' or ``composite'' that the Self assumes for its realization. On the other hand, a particular manifestation is not ``destroyed'' but only ``recomposed'' by being modified according to the states it passes through, be they heavenly or infernal, deriving into higher or lower forms. The key point is that the psyche is not a ``unit'', but only receives its ``unity'' as reflected by the transcendental Self.


\hfill

\end{sffamily}
\end{footnotesize}

